% =====================================================================
%  DRIFTGUARD — WTMC 2026 submission (long paper, 10 pages + 4 appendix).
%  Curated narrative build. Cover message = C3 (label-free self-repair);
%  the imperfect-oracle / implicit-regularization finding (C4) is the
%  strong secondary twist.
%
%  HOW THIS FILE IS ORGANISED
%  --------------------------
%  Each section lives in its own .tex and is pulled in via \input below.
%  We write them ONE AT A TIME; until a section file exists, keep its
%  \input line commented so the skeleton compiles on Overleaf.
%
%  Section map (agreed structure):
%    I    Introduction              introduction.tex
%    II   Related Work              related_work.tex
%    III  System Design             system_design.tex
%    IV   Experimental Evaluation   evaluation.tex
%    V    Operational Viability & Safety   safety.tex
%    VI   Oracle Quality: Multi-LLM Benchmark   oracle_benchmark.tex
%    VII  Discussion & Limitations  discussion.tex
%    VIII Conclusion & Future Work  conclusion.tex
%    App  Appendix                  appendix.tex
%
%  FIGURES: the source draft references ../runs/... ; on Overleaf upload
%  the PDFs into a figures/ folder and either keep that path or set
%  \graphicspath. A commented \graphicspath is provided below.
% =====================================================================

\documentclass[10pt,sigconf,letterpaper,nonacm]{acmart}

\AtBeginDocument{\providecommand\BibTeX{{Bib\TeX}}}

\settopmatter{printacmref=false}
\setcopyright{none}

\usepackage{subcaption}
\usepackage{paralist}
\usepackage{float}
\usepackage{multirow}
\usepackage{enumitem}
\usepackage{multicol}

\setlength{\textfloatsep}{6pt}   % space between floats and text
\setlength{\floatsep}{6pt}       % space between floats

% \graphicspath{{figures/}{../runs/al/_aggregate/}}

\begin{document}

\title[DRIFTGUARD: Label-Free Self-Repair for Phishing Detection]%
{DRIFTGUARD: Label-Free Self-Repair of Phishing Detectors under Drift\\
via LLM-Guided Active Learning}

\author{Simone Capriolo}
\affiliation{
 \institution{Austrian Institute of Technology\\Politecnico di Torino}
 \city{Vienna}
 \country{Austria}}
\email{simone.capriolo@ait.ac.at}

\author{Pedro Casas}
\affiliation{
  \institution{Austrian Institute of Technology}
  \city{Vienna}
  \country{Austria}}
\email{pedro.casas@ait.ac.at}

\author{Matteo Boffa}
\affiliation{
  \institution{Politecnico di Torino}
  \city{Turin}
  \country{Italy}}
\email{matteo.boffa@polito.it}

\renewcommand{\shortauthors}{Capriolo et al.}

\begin{abstract}
Machine-learning phishing detectors are trained on a snapshot of the threat
landscape and begin to decay the moment they ship: in our setting, a
leakage-clean detector scoring $98\%$ weighted $F_1$ in-distribution falls to
$\approx\!63\%$ on temporally fresh phishing. The standard remedy --- route
low-confidence samples to human analysts and retrain --- scales with the threat,
not with the budget. We present DRIFTGUARD, a closed-loop phishing detector that
repairs itself under drift with \emph{no human labels} and a small, predictable
LLM budget. DRIFTGUARD detects drift from the \emph{input} distribution rather
than from model confidence --- which a well-calibrated detector keeps high even
as it fails --- spends its budget on samples that are both uncertain and aligned
with the measured drift, and queries a large language model as an on-demand
oracle for gold labels derived from cleaned HTML, the opinions of two frozen
base detectors, and OSINT-derived domain intelligence. Retraining on the full
union of its seed corpus and every acquired label, DRIFTGUARD recovers nearly
$18$ points of weighted $F_1$ on fresh phishing ($0.63\!\to\!0.79$, five-seed
mean) using $\sim\!1{,}250$ selective queries at a cost of a few dollars, with no
measurable forgetting. Two design choices make this recovery trustworthy and
yield a counter-intuitive finding. First, a single cross-dataset statistical
gate sterilizes the feature space, rejecting both classic in-distribution
shortcuts and temporally leaky OSINT features under one rule, so every reported
gain reflects the phishing phenomenon rather than a dataset artifact. Second, the
oracle's imperfection is an asset, not a liability: a class-asymmetric verifier
that filters its label noise yields \emph{no} improvement on drift recovery and
consistently \emph{degrades} out-of-distribution robustness, because the filtered
noise was acting as an implicit regularizer --- so DRIFTGUARD keeps every paid
label. We further benchmark six LLM providers as oracles and stress-test a
label-free health monitor against live prompt injection, delimiting what
self-repair can and cannot guarantee in deployment.
\end{abstract}

\keywords{phishing detection; temporal concept drift; active learning;
large language models; LLM-as-oracle; data leakage; OSINT measurement;
implicit regularization; label-free monitoring}

\maketitle

% --- Sections (uncomment each as its file is written) ----------------
% \input{introduction}      % Section I
% \input{related_work}      % Section II
% % =====================================================================
%  DRIFTGUARD — Section III: System Design.
%  Narrative rewrite (NOT a clean-up of reboot/docs/System_Design.tex).
%  Facts/numbers sourced from the bible; prose written fresh, story-first.
%  Method only; all measurements live in Section IV (\ref{sec:eval}) and
%  the safety layer in Section V (\ref{sec:safety}).
%
%  FIGURE: needs an architecture diagram (figures/architecture.pdf).
%  Placeholder included below — produce the diagram before submission.
% =====================================================================

\section{System Design}
\label{sec:design}

A phishing detector that must survive drift faces a tension. It has to answer
in milliseconds at the edge, yet the knowledge it answers with goes stale in
weeks. DRIFTGUARD resolves this by refusing to do everything in one place. A
small, fast model gives the verdict; a separate, label-free signal decides
\emph{when} the model has falle n behind; and only then, on a tight budget, a
large language model is consulted to manufacture the few gold labels needed to
bring the model back. The design mirrors how a security operations centre
actually runs --- cheap automated triage in line, expensive expert judgement
on demand --- and turns the LLM from an always-on dependency into a periodic
recalibration tool (Figure~\ref{fig:arch}).

Two principles run through the whole system and are worth stating up front.
First, \emph{honesty before recovery}: before any adaptation, we sterilize the
feature space of signals that look predictive only because of how the data was
collected, so that the accuracy we later recover is real and not an artefact.
Second, \emph{nothing informative is wasted}: the signals we forbid the model
from using are not discarded but handed to the oracle, which can exploit
contextually what a statistical model can only exploit as leakage. The rest of
this section develops the three components in turn; the operational safety
layer that wraps them is deferred to Section~\ref{sec:safety}.

\begin{figure}[t]
  \centering
  % TODO: produce figures/architecture.pdf (edge model -> pool -> drift
  % detector -> selection -> LLM oracle -> full-union retrain -> snapshot,
  % with the safety layer auditing the oracle stream).
  \includegraphics[width=\columnwidth]{figures/architecture.pdf}
  \caption{DRIFTGUARD as a closed loop. The leakage-clean detector serves
  instant verdicts at the edge and quietly buffers the samples it is unsure
  about. A label-free monitor watches the \emph{input} distribution and, when
  it has shifted enough, triggers a round: a small batch of uncertain,
  drift-aligned samples is sent to the LLM oracle for gold labels, and the
  detector is retrained on the full union of its seed corpus and every label
  acquired so far. A safety layer (Section~\ref{sec:safety}) audits the oracle
  stream throughout.}
  \label{fig:arch}
\end{figure}

\subsection{A Leakage-Clean Detection Core}
\label{ssec:core}

The detector is a light meta-learner that sits on top of two existing,
pre-trained phishing classifiers used as frozen \emph{proposers}: a graph
neural network over the rendered DOM and a gradient-boosted model over HTML
features. Neither is retrained inside the loop; each contributes a phishing
probability and the two are combined with simple agreement features. On their
own, the proposers are brittle --- under temporal drift the DOM-structure
signal essentially collapses to chance --- which is precisely why we do not
trust them blindly and instead ground the verdict in evidence they do not see.

That evidence comes from a \emph{passive OSINT measurement layer}. For every
page we measure the lifecycle of its domain from free, key-less public sources:
its popularity (the Tranco list), its archival history (the Wayback Machine),
its certificates (Certificate Transparency logs), its registration record
(RDAP), and its DNS and network footprint. We deliberately restrict ourselves
to sources that carry no label leakage by construction --- popularity,
archival, registration, routing --- and exclude anything correlated with the
phishing ground truth, such as blocklists, from the model's feature space. The
result is a few dozen factual lifecycle measurements per page, alongside the
proposer opinions and a handful of URL and HTML features.

The methodological centrepiece is what happens next. Many of these features are
strongly predictive on the training corpus yet betray the model under drift,
and the reason is subtle: a feature like ``number of Wayback snapshots''
separates the classes almost perfectly on historical data not because it
distinguishes phishing from benign, but because the historical phishing domains
are now \emph{dead} and therefore unarchived --- it measures ``alive versus
dead'', a temporal accident, not the phenomenon we care about. We catch such
features with a single statistical rule: any feature whose single-feature
discriminative power is suspiciously high on the training set \emph{and}
inconsistent across datasets is rejected as leakage. The same rule, applied
uniformly, turns out to catch two failures of completely different nature --
the temporally leaky OSINT signals \emph{and} the proposers' own
near-perfect-but-non-transferable probabilities --- under one criterion. We
quantify exactly which features it removes, and what that does to
cross-distribution accuracy, in Section~\ref{ssec:eval-gate}.

Crucially, a rejected feature is removed from the \emph{model}, not from the
\emph{system}: the gated-out lifecycle signals are preserved and forwarded to
the oracle, where contextual reasoning can use them safely. On the sterilized
feature space we deploy a logistic-regression detector. It starts from a
high-precision, low-recall operating point --- the pages it calls phishing are
almost always phishing, but under drift it misses many --- and those misses are
exactly the region the self-repair loop is built to recover.

\subsection{An LLM Oracle for Gold Labels}
\label{ssec:oracle}

When the loop needs a label it cannot get for free, it asks a large language
model. The oracle's job is narrow but demanding: given everything we know about
a single page, return a verdict reliable enough to retrain on. We give it three
complementary views of the page. The first is the cleaned HTML --- stripped of
styling and binary blobs but keeping the forensic markers (hidden fields,
password inputs, form targets, inline scripts) a verdict actually turns on. The
second is the OSINT dossier, including the very lifecycle signals the leakage
gate removed from the model: presented as plain facts (``domain registered 207
days ago; never archived; certificate rotated every 90 days''), they let the
oracle reason about throwaway phishing infrastructure that the statistical
model was rightly forbidden from exploiting. The third, when useful, is the
proposers' opinions, framed as advice rather than truth.

We ask the oracle to reason in fixed steps and to commit to a decision rather
than hedge, and we require it to list both the indicators it found and the
checks that came up empty. This last requirement is not cosmetic: the two
parallel lists are what later lets a verifier check each claim against the page,
turning an opaque verdict into an auditable one (Section~\ref{sec:safety}). The
output is a single structured record, robust to malformed generations, with the
verdict drawn from a closed taxonomy so that downstream analysis can reason
about \emph{why} a page was flagged, not just whether. Which combination of
evidence makes the best oracle is itself an empirical question, settled in
Section~\ref{ssec:eval-oracle}.

\subsection{The Self-Repair Loop}
\label{ssec:loop}

The loop ties the detector and the oracle together and is where the
``self-repair'' claim is earned. It must answer three questions without ever
seeing a ground-truth label at runtime: when has the model fallen behind, which
handful of pages are worth a costly oracle call, and how to fold the answers
back in without forgetting what the model already knew.

\paragraph{When: watching the input, not the confidence.}
The instinctive drift signal --- the model losing confidence on new data --- is
the wrong one here, and for an instructive reason. A leakage-clean,
well-calibrated detector does \emph{not} get hesitant under drift; it stays
confident and is simply wrong, making confident mistakes on inputs it has never
seen. Any trigger that waits for confidence to collapse therefore waits
forever. DRIFTGUARD instead measures the shift in the \emph{input}
distribution: it compares the running distribution of the model's predictions
against a reference taken at deployment, and fires when that divergence grows.
We keep this timing mechanism deliberately simple --- in most runs, a scheduled
batching policy --- precisely so that the gains we report can be attributed to
\emph{which} samples we label rather than to a bespoke trigger. The operational
consequences of this choice, including the cost of reacting late, are measured
in Section~\ref{ssec:eval-delay}.

\paragraph{Which: spending the budget where the drift is.}
When a round fires, the detector quietly maintains a pool of the samples it was
uncertain about, and we must pick a small batch to label. Classical active
learning picks the most uncertain points; but ``uncertain'' is defined relative
to the current model, blind to \emph{where the data has actually moved}. Our
selector, \emph{drift-anchored} sampling, corrects this. It scores a candidate
by its boundary uncertainty, then up-weights it by how far it has drifted from
the training distribution along the dimensions that have drifted most:
\[
s(x) \;=\; \underbrace{\frac{1}{|p(x)-0.5|+\epsilon}}_{\text{uncertainty}}
\;\cdot\;
\Big(1 + \alpha \textstyle\sum_i \delta_i\,|x_i-\mu_i|\Big),
\]
where $\delta_i$ is the measured per-feature drift and $\mu_i$ the training
mean. The effect is to spend labels on pages that are both hard \emph{and}
explanatory of the observed shift, degrading gracefully to plain uncertainty
sampling when nothing has drifted. We benchmark it against the standard active-
learning strategies, and surface where it helps and where it trades off, in
Section~\ref{ssec:eval-recovery}; the single weight $\alpha$ is shown not to
matter in Section~\ref{ssec:eval-alpha}.

\paragraph{How: retraining without forgetting.}
Each round, the detector is refit from scratch on the \emph{full union} of its
original corpus and every label acquired so far. This is the blunt but reliable
guarantee against catastrophic forgetting: because all of the original evidence
is re-presented at every fit, recovery on the drifted distribution cannot come
at the price of the original one --- a claim we verify rather than assume
(Section~\ref{ssec:eval-recovery}). The feature scaler, fit once at deployment,
is frozen for the life of the run, since re-estimating it on a few dozen new
samples each round would induce far more coefficient drift than it removes.
After each retrain the uncertainty pool is re-scored under the new model, so
pages that have become easy graduate out and the pool stays aligned with the
model's \emph{current} blind spots.

\medskip\noindent
This is the whole loop: an honest baseline, a contextual oracle fed the very
signals the baseline could not safely use, and a budgeted, drift-aware
acquisition cycle that retrains without forgetting. What we have not yet
addressed is how to run it safely when the oracle is an LLM and no ground truth
is available to check it --- the subject of Section~\ref{sec:safety}.
     % Section III
% % =====================================================================
%  DRIFTGUARD — Section IV: Experimental Evaluation.
%  Narrative rewrite (NOT a clean-up of Experimental_Evaluation.tex).
%  Numbers are taken verbatim from the bible; prose is fresh, story-first.
%  Six tables kept in-body: data, leakage, xdataset, oracle, al-strategy,
%  cursor. The alpha-sensitivity table are deferred to
%  the appendix; only their takeaways appear here.
% =====================================================================

\section{Experimental Evaluation}
\label{sec:eval}

Our evaluation follows the life of the system and answers four questions in
order. Does the leakage gate buy us an \emph{honest} baseline, one whose
cross-distribution accuracy reflects phishing rather than collection artefacts
(Section~\ref{ssec:eval-gate})? Given that baseline, which evidence makes the
LLM the best possible oracle (Section~\ref{ssec:eval-oracle})? How much drift
does the self-repair loop actually recover, at what operational cost, and
without forgetting (Section~\ref{ssec:eval-recovery})? And finally, why is an
\emph{imperfect} oracle not merely tolerable but causally \emph{preferable} for
robustness (Section~\ref{sec:noise})? Unless noted, multi-seed results are
reported over five clean seeds as mean $\pm$ 95\% confidence interval;
single-stream operational analyses use the representative seed $2025$.

\subsection{Datasets and Protocol}
\label{ssec:eval-data}

We use three datasets, each with a fixed and distinct role
(Table~\ref{tab:data}). The detector is trained \emph{only} on KnowPhish (KP).
PhreshPhish (PP), collected in 2024--25, is a temporally disjoint corpus of
fresh phishing and is our principal drift test set. DeltaPhish/SpacePhish (DP)
is held out entirely as an adversarial out-of-distribution (OOD) probe and is
\emph{never} trained on. The two kept apart deliberately: KP$\to$PP measures
natural temporal drift, KP$\to$DP measures adversarial robustness, and conflating
them is exactly how prior work overstates generalisation.

For the self-repair loop, KP is split $80/20$ into a training fold (carried
unchanged into every retrain) and an in-distribution test (the forgetting
check). PP is replayed as a row-by-row \emph{stream} in chronological order ---
so the simulation sees drift as it happened --- with a disjoint held-out test
set the oracle never touches. Three invariants are enforced in code and matter
for every number that follows: DP labels never enter training, the PP test set
is isolated from both the stream and the labelling pool, and the drift detector
sees only the model's own output, never a ground-truth label.

\begin{table}[H]
\centering
\caption{The data and its operational slices ($\mathrm{seed}=2025$). Only
$\mathrm{KP_{train}}$ and accepted oracle labels are ever trained on; DP is a
pure OOD probe.}
\label{tab:data}
\small
\begin{tabular}{@{}llrl@{}}
\toprule
Slice & Source & $N$ & Role \\
\midrule
$\mathrm{KP_{train}}$ & KP 80\% & 7{,}221 & seed training set (frozen) \\
$\mathrm{KP_{test}}$  & KP 20\% & 1{,}806 & forgetting check \\
$\mathrm{PP_{stream}}$ & PP 80\% & 3{,}821 & drift stream (1{,}865 phish) \\
$\mathrm{PP_{test}}$  & PP 20\% & 956 & drift-recovery metric (held out) \\
$\mathrm{DP_{probe}}$ & DP full & 6{,}426 & OOD probe (85\% benign, eval only) \\
\bottomrule
\end{tabular}
\end{table}

\subsection{An Honest Baseline}
\label{ssec:eval-gate}

The leakage gate rejects $13$ of $39$ candidate features, and the rejections
tell a coherent story (Table~\ref{tab:leakage}). The archival and popularity
signals --- Wayback, Tranco, DNS --- are near-perfect on KP yet no better than a
coin flip on DP: as anticipated in Section~\ref{ssec:core}, they encode
``alive versus dead'', a property of \emph{when} the data was collected, not of
phishing. The same rule, with no special-casing, also rejects the proposers'
own headline probabilities (\texttt{max\_prob}, \texttt{prob\_sn}): brilliant on
KP, useless on fresh phishing. That a single statistical criterion catches a
shortcut in the base detectors and a temporal artefact in the OSINT layer ---
two failures of entirely different origin --- is the methodological payoff of
this stage. What survives is consistent across datasets: of the proposer
signals only \texttt{min\_prob} and \texttt{prob\_pg} generalise, joined by
registration-age features that are robust because domain age is genuinely
informative.

\begin{table}[H]
\centering
\caption{Representative gate decisions. AUC is the polarity-corrected
single-feature AUC per dataset. The same rule rejects temporally leaky OSINT
signals \emph{and} the proposers' non-transferable probabilities.}
\label{tab:leakage}
\small
\begin{tabular}{@{}lcccl@{}}
\toprule
Feature & $\mathrm{AUC}_{kp}$ & $\mathrm{AUC}_{dp}$ & $\mathrm{AUC}_{pp}$ & Verdict \\
\midrule
\texttt{wayback\_n\_snapshots}   & 0.996 & 0.504 & 0.914 & reject \\
\texttt{max\_prob}               & 0.989 & 0.802 & 0.504 & reject \\
\texttt{prob\_sn}                & 0.978 & 0.761 & 0.513 & reject \\
\texttt{tranco\_rank}            & 0.962 & 0.500 & 0.655 & reject \\
\midrule
\texttt{min\_prob}               & 0.982 & 0.978 & 0.858 & keep \\
\texttt{prob\_pg}                & 0.976 & 0.987 & 0.863 & keep \\
\texttt{rdap\_domain\_age\_days} & 0.884 & 0.503 & 0.859 & keep \\
\bottomrule
\end{tabular}
\end{table}

Sterilising the feature space costs almost nothing in-distribution --- all three
candidate estimators still exceed $98\%$ weighted $F_1$ on KP --- but changes the
picture out of distribution entirely (Table~\ref{tab:xdataset}). Ungated, the
meta-learner degenerates on the adversarial probe, calling almost everything
phishing ($F_{1,w}=7.1\%$); gated, that pathology is gone. Among the three
estimators, logistic regression maximises PP, the drifted distribution we care
about, and is the deployed detector. It begins where a good self-repair target
should: high phishing precision ($0.843$) and low recall ($0.372$). The pages it
calls phishing are almost always phishing; the ones it misses are exactly the
territory the loop must reclaim. We adopt its $\mathrm{PP}\,F_{1,w}=0.629$ as the
leakage-clean baseline. This is, by design, \emph{below} the stronger base
detector taken alone (PhishGraph reaches $0.748$ on PP) --- we have thrown away
the shortcuts that inflate it --- and the contribution is not a higher starting
point but a system that, after adaptation, climbs past the static ceiling to
$0.791$ while a frozen detector cannot move at all.

\begin{table}[H]
\centering
\caption{Weighted $F_1$ per split under the gated feature space. LogReg
maximises PP, the drift target, and is deployed. Ungated, the same pipeline
collapses to $7.1\%$ on the adversarial DP probe.}
\label{tab:xdataset}
\small
\begin{tabular}{@{}lccc@{}}
\toprule
Model & KP test & DP full & PP full \\
\midrule
LogReg & 0.981 & 0.514 & \textbf{0.629} \\
GB     & \textbf{0.990} & 0.481 & 0.592 \\
RF     & 0.989 & \textbf{0.535} & 0.592 \\
\bottomrule
\end{tabular}
\end{table}

\subsection{Seeing vs.\ Cheating: A Teacher--Student Gap}
\label{ssec:eval-gap}

The gate makes an asymmetric choice: it strips eleven lifecycle features ---
Tranco, Wayback, DNS, ASN --- from the detector, yet forwards them to the oracle
(Section~\ref{ssec:core}). Is that the right split, or does it needlessly blind
the detector? Three observations show the asymmetry is exactly right and, in the
process, reveal a genuine teacher--student gap that is a secondary message of
this paper.

First, the oracle truly leans on this evidence. Across $4{,}485$ cached verdicts,
$93\%$ of its stated reasoning invokes signals the detector cannot represent, and
two in three explicitly name Tranco or Wayback: the oracle is reasoning about a
domain's lifecycle, not only its page.

Second, those same signals are \emph{leakage} for a statistical model.
\texttt{wayback\_n\_snapshots} scores AUC $0.996$ on KP and $0.914$ on PP but
exactly chance ($0.504$) on the adversarial probe (Table~\ref{tab:leakage}): it
tracks ``alive versus dead'', which lines up with the label only because of when
each corpus was collected.

Third --- the crux --- if we hand the eleven features back to the detector, it
does not learn what the oracle knows; it learns to cheat
(Figure~\ref{fig:restoration}). The detector seizes on them eagerly: they take
$37\%$ of its weight and \texttt{wayback\_n\_snapshots} becomes its single most
important feature. But the gain is leakage, not understanding. In a controlled
replay on a fixed sample path, the restored detector reaches
$\mathrm{PP}\,F_{1,w}=0.899$ --- \emph{above} the $0.859$ the very same pipeline
attains with perfect ground-truth labels (and above the main loop's $0.873$
ceiling of Section~\ref{sec:noise}) --- a level only the temporal artefact can
buy, while KP saturates from $0.968$ to $0.996$. The lift lands precisely on the
time-correlated corpora and corresponds to nothing the model could generalise
with: on the adversarial probe the same feature is chance-level
(Table~\ref{tab:leakage}), so there is no shortcut to take.

\begin{figure}[t]
  \centering
  \includegraphics[width=\columnwidth]{figures/fig_restoration.pdf}
  \caption{The student can only cheat with the teacher's features. Handed the
  eleven gated-out OSINT signals, the statistical detector reaches a PP score
  (restored, $0.899$) \emph{beyond} what the same pipeline achieves with perfect
  ground-truth labels (the $0.859$ ceiling) --- impossible without exploiting the
  temporal leakage the gate removed --- and saturates KP. The gain is confined to
  the time-correlated corpora; on the adversarial probe the same feature is
  chance-level. Absolute values are from a controlled fixed-path replay and are
  not comparable to the main loop (Table~\ref{tab:al-strategy}); the
  gated-vs-restored contrast and the ceiling crossing are the point.}
  \label{fig:restoration}
\end{figure}

The picture is now coherent. Detector and oracle see the same signal but can use
it in qualitatively different ways. A statistical model can use it only as a
correlation, and that correlation \emph{is} the temporal leakage the gate
removed; the oracle reads it in context --- ``recently registered, never
archived, absent from Tranco'' together spell throwaway infrastructure --- and so
generalises. The teacher--student gap is therefore not that the student lacks
features: it is that the only way the student \emph{can} use them is the one the
gate rightly forbids. The student does not fail to \emph{see}; it can only
\emph{cheat} --- which is exactly why the same lifecycle evidence that would
corrupt the detector is what makes the oracle worth querying.

\subsection{The Right Evidence for the Oracle}
\label{ssec:eval-oracle}

Before wiring the oracle into the loop we ask what to show it, on the cases that
matter: the $100$ most diagnostic detector errors on PP --- the $50$ phishing it
was most sure were benign and the $50$ benign it was most sure were phishing ---
so that every sample carries a detector probability that is confidently
\emph{wrong}. The result is counter-intuitive: \textbf{handing the oracle the
detector's probabilities makes it worse} (Table~\ref{tab:oracle}). Removing them
lifts weighted $F_1$ by $22$ points and raises recovery of missed phishing from
$24/50$ to $40/50$. The mechanism is plain once stated --- on these samples the
probability is a confidently wrong prior, and the model partially anchors to it
--- and it is, read positively, a robustness property: the factual OSINT dossier
gives the oracle independent ground to stand on.

\begin{table}[H]
\centering
\caption{Oracle ablation on $100$ boundary samples ($50$ FN, $50$ FP). FN/FP
recovery is the fraction of the detector's errors the oracle corrects.
Probability-blind variants (C, D) dominate on these confidently-wrong cases.}
\label{tab:oracle}
\small
\begin{tabular}{@{}lccccc@{}}
\toprule
Variant & Acc & $F_{1,w}$ & $R_{\text{phish}}$ & FN rec. & FP rec. \\
\midrule
D \texttt{html\_osint}        & \textbf{0.910} & \textbf{0.909} & 0.820 & 41/50 & 50/50 \\
C \texttt{html\_only}         & 0.900 & 0.899 & 0.800 & 40/50 & 50/50 \\
E \texttt{poc\_2k}            & 0.820 & 0.814 & 0.640 & 32/50 & 50/50 \\
B \texttt{html\_probs\_osint} & 0.770 & 0.761 & 0.580 & 28/50 & 48/50 \\
A \texttt{html\_probs}        & 0.690 & 0.676 & 0.480 & 24/50 & 45/50 \\
\bottomrule
\end{tabular}
\end{table}

There is a twist that governs deployment. This test set is, by construction,
adversarial to the probabilities, which is why the probability-blind variants
win on it. The samples the loop actually queries are different: they sit near
the genuine decision boundary, where a probability encodes ``uncertain'' rather
than ``confidently wrong'', and there the relationship inverts and the full
multimodal variant becomes the stronger oracle. The lesson is that the best
oracle context is distribution-dependent --- probabilities help when they encode
real uncertainty and hurt when they encode confident error --- and we accordingly
adopt the full multimodal variant inside the loop, where its samples live.

\subsection{Recovery, Forgetting, and Selection}
\label{ssec:eval-recovery}

This is the main result (Table~\ref{tab:al-strategy}). Over five clean seeds the
self-repair loop lifts $\mathrm{PP_{test}}$ weighted $F_1$ from $0.623$ to
$\mathbf{0.791 \pm 0.013}$ --- about $+17$ points --- and it does so in the right
way: the gain is driven entirely by phishing recall ($0.37 \to 0.61$) while
precision is \emph{preserved} ($\approx 0.96$). The loop recovers missed
phishing without sacrificing the high-precision operating point that made the
detector trustworthy. And it does so without forgetting: across all $40$ runs the
in-distribution KP score never leaves a one-point band ($0.972$--$0.982$), the
empirical signature of full-union retraining. At roughly one to two dollars of
tokens per run, this is the headline the paper promises.

We compare our \emph{drift-anchored} selector (Section~\ref{sec:design}) against
the standard active-learning acquisition rules, each of which decides which pooled
samples to send to the oracle. \emph{Margin} is classical uncertainty sampling:
the samples whose predicted probability sits closest to the $0.5$ boundary.
\emph{QBC} (query-by-committee) picks the samples on which a committee of models
most disagrees. \emph{Core-set} ignores uncertainty and instead chooses a diverse
subset that geometrically covers the feature space (a $k$-center rule).
\emph{BADGE} fuses the two, clustering gradient embeddings that are large when the
model is uncertain and spread out when samples are diverse. \emph{Hybrid} is a
fixed blend of margin, core-set diversity, and committee disagreement.
\emph{Random} draws uniformly from the pool --- the budget-matched control that
isolates how much the \emph{choice} of samples matters beyond their number.
Finally, the \emph{ungated stream} row removes the uncertainty pool entirely and
labels the raw delegate stream, isolating the pool's contribution.

\begin{table}[H]
\centering
\caption{Selection strategies over five seeds (periodic trigger, multimodal
oracle). PP is weighted $F_1$; DP is macro $F_1$ (the probe is $85\%$ benign);
``oracle acc'' is the agreement of accepted labels with ground truth. Round-0
baselines: $\mathrm{PP}=0.623$, $\mathrm{DP}=0.438$. The informed strategies are
indistinguishable on PP but separate sharply on the OOD probe.}
\label{tab:al-strategy}
\small
\begin{tabular}{@{}lccc@{}}
\toprule
Strategy & $\mathrm{PP}\,F_{1,w}$ & $\mathrm{DP}\,F_{1,\mathrm{mac}}$ & oracle acc \\
\midrule
drift-anchored (ours) & $\mathbf{0.791 \pm 0.013}$ & $0.547 \pm 0.023$ & 0.81 \\
margin     & $0.780 \pm 0.018$ & $0.706 \pm 0.015$ & 0.70 \\
qbc        & $0.779 \pm 0.019$ & $0.714 \pm 0.004$ & 0.70 \\
hybrid     & $0.779 \pm 0.019$ & $0.714 \pm 0.008$ & 0.71 \\
core-set   & $0.778 \pm 0.017$ & $0.687 \pm 0.012$ & 0.73 \\
badge      & $0.776 \pm 0.021$ & $0.708 \pm 0.012$ & 0.72 \\
random     & $0.762 \pm 0.016$ & $0.634 \pm 0.016$ & 0.87 \\
\midrule
\emph{ungated stream} & $0.766 \pm 0.011$ & $0.598 \pm 0.025$ & 0.84 \\
\bottomrule
\end{tabular}
\end{table}

The selection results then say something we did not expect. On PP, the informed
strategies are statistically indistinguishable --- they cluster between $0.776$
and $0.791$ --- so the choice of acquisition rule is \emph{not} where the
in-distribution gain comes from. Where it matters is the adversarial probe, and
there a clean trade-off appears: drift-anchored attains the best PP but the
\emph{worst} DP, while the more diffuse strategies give up about one PP point to
buy $14$--$17$ points of DP. Concentrating labels along the measured drift
sharpens a PP-specific boundary; spreading them out generalises to the
orthogonal adversarial shift. Two controls pin the cause down. Randomising
\emph{when} the loop fires, with the selector fixed, leaves PP essentially
unchanged ($0.798$ vs $0.800$), confirming that timing is not the lever; removing
the uncertainty pool and labelling the raw stream costs $2.5$ PP points,
confirming that the pool is. The trade-off itself is real and we revisit its
cause next: notably, the cleanest-label strategy (drift-anchored, oracle accuracy
$0.81$) is the OOD-weakest, a clue that the label \emph{noise} is doing useful
work.

\paragraph{Robustness to the one knob.}
\label{ssec:eval-alpha}
Drift-anchored carries a single weight, $\alpha$, on its drift term. Sweeping it
across five values shows the behaviour is a \emph{step}, not a tuned magnitude:
at $\alpha=0$ the selector is exactly margin sampling, turning the term on at any
positive weight reshuffles most of the selected set and produces the PP/DP
profile above, and increasing $\alpha$ further changes essentially nothing. The
specific value is immaterial --- it is the \emph{presence} of the drift term, not
its tuning, that defines the method (appendix).

\subsection{Choosing the Trigger: the Cost of Delay}
\label{ssec:eval-delay}

Timing is not our contribution, but the loop must fire on \emph{some} signal, and
the choice carries a measurable operational cost: between the onset of drift and
the first repair, the stale model keeps misclassifying in production. Replaying
the stream under whichever snapshot was live at each row turns that cost into a
number --- phishing missed before repair, and the time-integrated error
$\int(1{-}F_{1,w})$ --- and lets us compare triggers fairly
(Table~\ref{tab:al-delay}).

\begin{table}[H]
\centering
\caption{Cost of delay on $\mathrm{PP_{stream}}$ ($3{,}821$ rows, $1{,}865$
phishing; seed $2025$), replayed under the model snapshot live at each row.
``missed'' is phishing missed in production before repair; $\int(1{-}F_{1,w})$ is
the time-integrated operational error; ``probe'' counts oracle calls spent only
to \emph{measure} agreement, yielding no training label. We show the two
deployable label-free triggers, the output-gated reference, and two
representatives of the agreement-based (LAR) family; the full sweep is in the
appendix.}
\label{tab:al-delay}
\small
\begin{tabular}{@{}lccccc@{}}
\toprule
Trigger & 1st & \#fires & missed & $\int(1{-}F_{1,w})$ & probe \\
\midrule
\textsc{kl-only}            & \textbf{99}  & 30 & \textbf{828}  & \textbf{830.6}  & 0 \\
\textsc{periodic} (ours)    & 149 & 25 & 866  & 875.0  & 0 \\
\textsc{or} (output-gated)  & 429 & 2  & 1178 & 1217.4 & 0 \\
\midrule
LAR (best: random, $w{=}100$) & 299 & 17 & 1070 & 1103.3 & 380 \\
LAR (budget, $w{=}200$)       & 599 & 2  & 1144 & 1176.4 & 190 \\
\bottomrule
\end{tabular}
\end{table}

Two label-free triggers are deployable. \textsc{kl-only} watches the input
distribution and fires the moment it shifts: it reacts earliest (row $99$),
misses the fewest phishing, and incurs the lowest integrated error, at the cost
of firing often. \textsc{periodic} fires on a fixed schedule: it reacts a little
later but spends a \emph{predictable} budget, which is why we adopt it as the
canonical trigger --- a sensible default for a proof-of-concept, and within a
hair of \textsc{kl-only} on every metric ($0.74$ vs $0.75$ stream-wide weighted
$F_1$). The \textsc{or} trigger, which waits for the model's \emph{confidence} to
collapse, is the cautionary reference: on a calibrated detector that confidence
never collapses, so it fires twice in the entire stream and never leaves the
baseline.

We also asked whether the oracle itself could time the loop. The natural idea is
to periodically \emph{probe} a few samples with the LLM and fire when the model's
agreement with it drops --- a model-versus-oracle disagreement signal we call the
LLM-agreement rate (LAR), which we tried in several variants (random or
drift-guided probing, relative-drop or absolute-threshold rules, two window
sizes). It is intuitive but loses on every axis (Table~\ref{tab:al-delay}, lower
block): the best variant reacts $200$ rows later than \textsc{kl-only}, misses
over $240$ more phishing, and --- worse --- burns several hundred extra oracle
calls purely to measure agreement, calls that produce no label to train on. The
reason is structural: an agreement trigger needs the \emph{oracle} to witness the
drift, whereas the same shift is already visible, for free, in the model's own
prediction histogram. Watching the input is both the better signal and the
cheaper one, and it is what we deploy.

\subsection{Why an Imperfect Oracle Wins}
\label{sec:noise}

The trade-off above hinted that the oracle's mistakes are not pure cost. We now
show they are causally beneficial for robustness, with two experiments.

First, the ground-truth ceiling. Replacing the LLM with perfect labels,
everything else fixed, bounds what relabelling could ever buy. On PP it is a true
ceiling of $0.873$, of which our $\approx\!\$1$ oracle already captures two
thirds. But on the adversarial probe the perfect oracle is not a ceiling --- it
is a \emph{floor}: the ground-truth-trained model regresses to a macro $F_1$ of
$0.337$, \emph{below} the untrained baseline of $0.438$ and far below the LLM
run's $0.547$. Perfect labels sharpen the boundary so precisely onto PP's
phishing region that the $85\%$-benign probe gets swept up as phishing. If
perfect labels hurt OOD robustness and imperfect ones help it, the \emph{noise}
is doing the work.

Second, we isolate that claim. A naive LLM-versus-ground-truth comparison
confounds two things --- the labels differ and the selected samples can differ
--- so we hold the queried samples fixed and vary only label quality, flipping an
increasing fraction $\varepsilon$ of the oracle's mistakes back to truth and
refitting (Table~\ref{tab:cursor}). The trend is monotone and clean: as the
labels get \emph{better}, PP rises ($+9$ points) while every OOD metric falls
--- macro $F_1$, PR-AUC, phishing $F_1$ all collapse --- and the in-distribution
score stays flat, isolating the effect to the adversarial axis. At
$\varepsilon=1$ the cursor lands exactly on the independent ground-truth ceiling,
confirming the two methods agree.

\begin{table}[H]
\centering
\caption{Holding the queried samples fixed and correcting the oracle's errors
toward ground truth ($\varepsilon\!:\!0\to1$, margin selector, seed $2025$). As
labels improve, PP rises and the OOD probe falls monotonically; KP is unaffected.}
\label{tab:cursor}
\small
\begin{tabular}{@{}lccccc@{}}
\toprule
$\varepsilon$ & $\mathrm{PP}\,F_{1,w}$ & $\mathrm{PP}\,R_{\text{ph}}$ & $\mathrm{DP}\,F_{1,\mathrm{mac}}$ & $\mathrm{DP}\,\text{PR-AUC}$ & $\mathrm{KP}\,F_{1,w}$ \\
\midrule
0.00 & 0.775 & 0.589 & \textbf{0.683} & 0.753 & 0.973 \\
0.25 & 0.803 & 0.644 & 0.661 & 0.733 & 0.977 \\
0.50 & 0.834 & 0.705 & 0.613 & 0.709 & 0.978 \\
0.75 & 0.855 & 0.751 & 0.522 & 0.663 & 0.979 \\
1.00 & \textbf{0.866} & 0.812 & \textbf{0.350} & 0.568 & 0.978 \\
\bottomrule
\end{tabular}
\end{table}

The reading is direct: the oracle's $\approx\!29\%$ error rate acts as label
noise that keeps the boundary from over-committing to the PP-specific direction,
and so preserves robustness to an orthogonal adversarial shift; remove the noise
and the probe falls below its own untrained baseline. This reframes oracle
imperfection from a limitation into a mechanism --- in an adversarial deployment a
moderate oracle error rate is a feature, not a defect. It also raises a sharp
operational question for any deployment --- whether one should \emph{clean} the
oracle's labels before training on them --- which Section~\ref{sec:safety} takes
up, with a more nuanced answer than the noise result alone would suggest.
        % Section IV
% % =====================================================================
%  DRIFTGUARD — Section V: Operational Viability and Safety.
%  Narrative rewrite. The verifier is presented in a BALANCED way: it is a
%  genuinely good label-quality instrument; whether to use it to gate
%  retraining is a budget/horizon-dependent trade-off, not a flat negative.
%  Figures: fig_verifier_overview.pdf, fig_injection_live.pdf.
%  Verifier ablation exact numbers + health-signal thresholds -> appendix.
% =====================================================================

\section{Operational Viability and Safety}
\label{sec:safety}

Recovering drift is one thing; running the loop safely is another. In
deployment the oracle is an LLM, its input is attacker-controlled page content,
and there is no ground truth at runtime to check it. This section builds the
safety layer that makes the loop deployable anyway, and is candid about its
limits. A class-asymmetric verifier scores every oracle label without an LLM
call; we show it is a good label-quality instrument, then ask -- and answer
with nuance -- whether it should \emph{gate} retraining. In addition, a label-free
health monitor is the loop's single circuit breaker, validated on oracles that
failed for real --- and we are explicit about the one failure mode it provably
cannot see, which we carry forward to future work.

\subsection{Scoring Oracle Labels without Ground Truth}
\label{ssec:verifier}

The oracle's two verdict classes are not equally reliable: on the AL stream a
phishing verdict is $\approx\!99\%$ precise, but a benign verdict only
$\approx\!59\%$ --- clean-looking scam pages the LLM excuses. The verifier is
built around that asymmetry. For each label it produces a soft trust score in
$[0,1]$ from three tiers: the model's own confidence, an evidence check that
dispatches each claimed indicator to one of a dozen HTML/URL/OSINT predicates
and asks whether it is actually present, and an OSINT-coherence tier that
\emph{reinforces} a phishing verdict but \emph{challenges} a benign one. For example a benign
call on a domain registered days ago, never archived and absent from Tranco is
penalised. The exact predicates and weights are in the appendix; two properties
matter here. It never needs an LLM call, and --- by design --- it never discards a
label on its own: it only measures how well a verdict is grounded in checkable
evidence.

As a measurement, it works well (Figure~\ref{fig:verifier}). Restricting the
canonical $1{,}247$-label stream to the labels the verifier trusts
($\tau\!\geq\!0.5$, $n{=}1{,}135$) raises oracle accuracy from $0.83$ to $0.89$
and, on the weak phishing class the loop most needs, oracle recall from $0.57$ to
$0.68$. The $112$ labels it distrusts are $72\%$ wrong, and essentially all of
them are the same failure: a genuinely phishing page the oracle excused as
benign. The verifier is, in short, a faithful label-quality instrument. The
question is what to \emph{do} with that measurement.

\subsection{To Gate, or to Keep Every Label?}
\label{ssec:verifier-gate}

The obvious move is to act on the trust score at training time --- softly, as a
per-label weight, or hard, dropping every label below $0.5$. We test both against
the no-verifier baseline in a \emph{controlled} ablation that replays the
identical $1{,}247$-label trajectory, so the only variable is the weight a label
receives (Figure~\ref{fig:verifier}, lines).

The result is a genuine trade-off, and reading it correctly matters. In
distribution, cleaning the labels costs essentially nothing: the full verifier
lands at $\mathrm{PP}\,F_{1,w}=0.799$ and the hard filter at $0.804$, against the
$0.800$ baseline, with KP flat throughout --- so feeding the model a demonstrably
cleaner label stream neither helps nor hurts what it does on the drifted
distribution. The one place gating leaves a mark is out of distribution, and it
is the mark Section~\ref{sec:noise} predicts: the labels the verifier removes are
precisely the imperfect ones that were acting as a regularizer, so dropping them
sharpens the boundary back toward PP and costs the adversarial probe $4$--$8$
points of macro-$F_1$, the more aggressively the more it filters. A second,
independent instrument thus reproduces the cursor of Section~\ref{sec:noise}: the
``noisy'' labels are not defects to be filtered, they are the regularizer.

We therefore keep \emph{every} label at uniform weight in DRIFTGUARD, preserving
the out-of-distribution robustness of Section~\ref{sec:noise}. But this is a
decision about a \emph{regime}, not a verdict that label hygiene is worthless.
The reason hygiene buys nothing measurable here is that the $7{,}221$-sample seed
corpus dominates the retrain loss, so re-weighting a thousand oracle labels
barely moves the boundary; the regularizing noise, by contrast, is felt on the
sensitive adversarial axis. In a long-running deployment the balance shifts: as
acquired labels accumulate and eventually rival the seed, and as label errors
compound across many retrains, a verifier that lifts label accuracy from $0.83$
to $0.89$ becomes an asset. The honest reading
is that the verifier's value is real and grows with scale and horizon, while its
small-budget out-of-distribution cost is simply the flip side of the very noise
that helps today --- and that, at our budget, keeping the labels is the right call.

\begin{figure}[t]
  \centering
  \includegraphics[width=\columnwidth]{figures/fig_verifier_overview.pdf}
  \caption{The verifier works --- which is what makes the gating question
  interesting. Bars: across modes, restricting to the labels the verifier trusts
  raises oracle accuracy ($0.83\!\to\!0.89$) and phishing recall
  ($0.57\!\to\!0.68$). Lines: feeding that cleaner stream to the retrain leaves
  in-distribution PP flat (even marginally higher for the hard filter) but costs
  out-of-distribution robustness (DP), exactly the regularization of
  Section~\ref{sec:noise}. DRIFTGUARD keeps every label at uniform weight; the
  hygiene the verifier offers is a budget- and horizon-dependent asset, not a
  one-shot $F_1$ boost.}
  \label{fig:verifier}
\end{figure}

\subsection{The Verifier as a Cross-Model Auditor}
\label{ssec:auditor}

If the verifier should not gate individual labels, its real, regime-independent
job is \emph{observability}. Because it scores the \emph{grounding} of a verdict
rather than its outcome, and does so without ground truth, it yields per-model
behavioural metrics that accuracy alone cannot: how often a provider's claims are
refutable (a hallucination proxy), how well its confidence tracks correctness,
and what reasoning style it favours. This makes it the natural instrument for
choosing an oracle and watching it drift across models --- exactly how it is used
in Section~\ref{sec:oracle-quality} --- and it feeds the health monitor one
round-level trust aggregate, never a per-label decision. That benchmark also
closes the present argument from the other side: it surfaces two robust paradoxes
--- the verifier's highest-trust labels are \emph{less} accurate than its
second-highest, and the samples it flags as hallucinating are \emph{more}
accurate than clean ones --- which independently confirm that neither trust nor
hallucination is safe as a per-label gate.

\subsection{The Label-Free Health Monitor}
\label{ssec:safety-layer}

The verifier measures but never acts; the single actuator is a health monitor
that detects a structurally broken oracle from the output stream alone, with no
ground truth. An oracle can degrade silently --- a wrong model behind a gateway, a
rate-limited provider returning malformed JSON, a page attacking the prompt ---
and quietly poison the retrain. Each round the monitor computes a few label-free
signals over the returned verdicts: the \emph{parse rate} (schema-valid
fraction), the \emph{class skew}, the \emph{yield}, and optionally the verifier's
round-mean trust. The parse rate is the primary discriminator, and deliberately a
parse-layer rather than a network-layer check: a broken oracle can return
HTTP~200 with an unparseable body, whereas a healthy one essentially always emits
schema-conforming output, so its collapse has no benign explanation. Skew, by
contrast, is \emph{not} treated as breakage on its own --- a drift-anchored batch
is legitimately near-mono-class early in the stream --- so its critical threshold
is set at true mono-class and no single signal can halt the run. The monitor
stops only on corroborated, sustained failure (two consecutive multi-signal
critical rounds, or three consecutive rounds of degraded parse rate); it is a
pure observer that never edits a label, so it cannot bias learning.

The clearest evidence that this is the right design is a case where it fired for
real (Figure~\ref{fig:llama-halt}). Swapping the oracle for a
\textsc{Llama-3.3-70B} endpoint, the loop runs healthy for fifteen rounds and
then begins to degrade: the provider's schema-valid fraction erodes round by
round until three consecutive rounds fall below the parse-warning line, tripping
the sustained-parse rule and halting the loop at round~18. The decisive detail is
what the monitor does \emph{not} see. It reads only the parse-ok rate; it is
blind to the test-time $F_{1,w}$ of the model being retrained --- which, on this
oracle, is quietly sliding the whole time. Acting on the structural signal alone
is nonetheless enough: the halt stops the loop before the degrading labels can do
lasting harm, and the monitored run ends strictly above where the unmonitored run
would have driven it. The same parse-layer logic catches the two other ways the
oracle failed in our benchmark, without a single false alarm: a \textsc{Grok-4}
endpoint that was dead on arrival ($0/20$ usable labels), and a \textsc{Kimi-K2.6}
whose calls all returned HTTP~200 yet were $68\%$ unparseable --- which a network
liveness check would have waved through and the parse rate catches at once. Across
the four healthy providers it never tripped in $25$ rounds each; its full
cross-model validation is reported with the multi-LLM benchmark
(Section~\ref{sec:oracle-quality}).

\begin{figure}[t]
  \centering
  \includegraphics[width=\columnwidth]{figures/fig_llama_halt.pdf}
  \caption{The health monitor halting a real failing oracle. With a
  \textsc{Llama-3.3-70B} oracle the schema-valid (parse-ok) rate --- the only
  signal the monitor reads --- erodes after round~15 and trips the sustained-parse
  rule (three consecutive rounds below $0.95$), halting the loop at round~18. The
  monitor is \emph{blind} to the model's test $F_{1,w}$ (dotted, right axis),
  which degrades throughout; acting on the structural signal alone is enough to
  stop the loop before the broken oracle does lasting harm.}
  \label{fig:llama-halt}
\end{figure}

The honest limit follows directly from how the monitor works: it is a
\emph{structural-breakage} detector --- malformed output, sustained parse failure,
gross class collapse --- not a \emph{label-correctness} detector. An oracle that
stays perfectly schema-conforming while a prompt-injected page coerces it into a
confident, well-formed \emph{wrong} verdict produces no signal a threshold can
catch, and no threshold change closes that gap. The natural defence is the
verifier's OSINT-coherence tier, which by construction challenges a benign verdict
on a freshly registered, unarchived domain that a label-flip cannot also forge;
turning that observability into a safe runtime signal, without re-introducing the
gating pathology of Section~\ref{ssec:verifier-gate}, is a problem we take up
explicitly in Section~\ref{sec:conclusion}.

\medskip\noindent
The two components form the operational safety layer with a clean division of
labour: the verifier \emph{measures} (per-label grounding and a round-level trust
signal, never gating), the monitor \emph{acts} (a sustained, corroborated halt,
never on a single label). Together they let the loop run without runtime ground
truth --- keeping every paid label, and with it the implicit-regularization
benefit of Section~\ref{sec:noise}, auditing the oracle continuously, and
stopping itself before a structurally broken oracle can poison the model. The
boundary of that guarantee is explicit: it defends against an oracle that
\emph{breaks}, not against one quietly coerced into confident, well-formed lies
--- the open problem we carry to Section~\ref{sec:conclusion}.
            % Section V
% \input{oracle_benchmark}  % Section VI
% \input{discussion}        % Section VII (Discussion & Limitations)
% \input{conclusion}        % Section VIII (Conclusion & Future Work)
% % =====================================================================
%  DRIFTGUARD — Appendix: the multi-LLM oracle benchmark, in full.
%  Narrative appendix (paper voice), self-contained: benchmark design,
%  headline accuracy + trust calibration, the two paradoxes, reasoning
%  styles, CoT case studies, the anomaly catalogue + dataset noise floor,
%  the Llama true-positive halt, and the matched-budget cross-check.
%  Sourced from reboot/docs/Oracle_Benchmark.tex + Appendix.tex (facts).
%  Requires: booktabs, multirow, enumitem. Defines \label{app:bench}.
%  References fig:llama-halt (defined in safety.tex) and the safety-layer
%  labels ssec:verifier, ssec:verifier-gate, ssec:auditor, sec:noise.
% =====================================================================

\appendix

\section{The Multi-LLM Oracle Benchmark}
\label{app:bench}

The main paper fixes one oracle and shows the loop recovers with it. Two
deployment questions remain, and they are the kind a reviewer is right to press:
\emph{which} LLM should serve as the oracle, and does the safety machinery of
Section~\ref{sec:safety} --- the verifier and the health monitor --- generalise
beyond the single model it was calibrated on? We answer both by running the
\emph{identical} canonical loop across a panel of providers and then turning the
verifier back on their outputs as a label-free auditor. The short version is
counter-intuitive on every axis: the choice of model matters far less than the
structure of the prompt, the verifier's own confidence signals must be read
against the grain, and the providers that fail do so in exactly the ways the
health monitor is built to catch. This appendix gives the full evidence.

\subsection{Benchmark design}
\label{app:bench-design}

Every cell runs the same canonical configuration --- drift-anchored selection,
prompt~B, periodic trigger, verifier off in-loop (audited retrospectively),
health monitor on --- over the full PhreshPhish stream at a budget-friendly
$K{=}20$ ($\approx 25$ rounds, $\approx 500$ calls per provider,
$\approx$\$15--30 for the panel). Only the provider varies, so any difference is
attributable to the model and not to the pipeline around it. Eight providers were
attempted; two returned too few usable labels to audit and became the
broken-oracle case studies of Appendix~\ref{app:bench-llama}, leaving six audited.
Because $K{=}20$ sits below the $K{=}50$ of the main results, the benchmark is
designed to measure provider \emph{behaviour} --- label accuracy, trust
calibration, hallucination, reasoning style --- rather than absolute recovery, so
the accuracy we report is oracle accuracy on the AL-queried boundary samples, the
cleanest quantity that is comparable across models.

\subsection{Accuracy and trust calibration}
\label{app:bench-acc}

Table~\ref{tab:bench-acc} reports accuracy and trust calibration over the full
$25$-round budget, and the first thing to read in it is a trap. Llama-3.3-70B tops
the raw accuracy column, but not because it is the best oracle: it \emph{halts}
at round~18 (Appendix~\ref{app:bench-llama}) and therefore never faces the hard
late rounds --- the Oct~22--24 phishing cluster --- that drag down every other
model's average. On a rounds-matched comparison, recomputing each provider over
the same first $18$ rounds Llama completed, Llama falls to third ($0.87$) while
Mistral-Large-3 ($0.90$) and DeepSeek-V4-Flash ($0.87$) lead, and every full-$25$
model gains $3$--$6$ points once those late rounds are excluded.

One reading survives both framings, and it is the result that matters: the top of
the ranking is \emph{not} the frontier closed-source tier. Open-weights and
European models lead at matched budget, the GPT-5 family sits mid-pack, and the
gaps are small --- $8.9$ points best-to-worst on the raw column, and far less once
budgets are matched. The decomposition says why. Phishing accuracy is uniformly
high across all six providers ($0.95$--$1.00$); the \emph{entire} cross-model
variance lives on the benign class ($0.73$--$0.85$) --- the same class asymmetry
the verifier of Section~\ref{ssec:verifier} is built around. The practical
implication is that the structured prompt, not the model, governs oracle quality,
which is exactly what makes the loop portable across providers. The matched-budget
end-to-end cross-check of Appendix~\ref{app:bench-smallbench} corroborates it from
the other direction, with providers clustering within $\approx 4$ points.

Trust calibration is the column to watch for safety. The Pearson $\rho$ between
verifier trust and correctness ranges from $0.43$ (GPT-5.1) down to $0.16$
(Mistral): a high-accuracy oracle whose trust is \emph{decoupled} from correctness
is the dangerous deployment case, because the verifier then reports confidence
everywhere and flags none of the errors. This is the first hint that trust is an
observability signal, not a gate --- a point the next two paradoxes settle.

\begin{table}[h]
\centering
\caption{Oracle accuracy on AL-queried samples and trust calibration, over the
full $25$-round budget; $n$ is the number of usable (queried-and-parsed) labels.
Phishing accuracy is uniformly high; cross-model variance is confined to the
benign class. $^{\dagger}$Llama halted at round~18
(Appendix~\ref{app:bench-llama}), so its $n$ --- and hence its full-run accuracy
--- covers an earlier, easier subset and is \emph{not} budget-comparable; on a
rounds-matched basis it is mid-pack (see text).}
\label{tab:bench-acc}
\small
\begin{tabular}{@{}lrcccc@{}}
\toprule
Provider & $n$ & Acc & Acc$_{\text{ph}}$ & Acc$_{\text{bn}}$ & $\rho_{\tau,c}$ \\
\midrule
Llama-3.3-70B$^{\dagger}$ & 323 & 0.864 & 1.000 & 0.854 & 0.38 \\
Mistral-Large-3   & 490 & 0.859 & 0.972 & 0.827 & 0.16 \\
GPT-5.4-mini      & 481 & 0.823 & 0.951 & 0.780 & 0.31 \\
GPT-5.4           & 482 & 0.813 & 0.982 & 0.764 & 0.37 \\
DeepSeek-V4-Flash & 479 & 0.808 & 0.979 & 0.764 & 0.30 \\
GPT-5.1           & 484 & 0.775 & 1.000 & 0.730 & \textbf{0.43} \\
\bottomrule
\end{tabular}
\end{table}

\subsection{Two paradoxes: why verifier trust is observability, not a gate}
\label{app:bench-paradox}

The cross-model audit exposes two robust paradoxes that, taken together, replicate
the gating negative result of Section~\ref{ssec:verifier-gate} on an entirely
different (unstructured) benchmark --- and explain why DRIFTGUARD uses the
verifier as an auditor (Section~\ref{ssec:auditor}) rather than a per-label
filter.

\paragraph{The hallucination paradox.} Samples on which a provider hallucinates
--- the verifier refutes at least one claim --- are \emph{more} accurate than
clean ones, for all six providers (Table~\ref{tab:bench-halluc}; Fisher exact
odds ratios $1.8$--$4.7$, $p<0.001$ for five of six, Llama the only borderline
case at $p=0.10$). The mechanism is not fabrication but verbosity: phishing pages
afford many concrete indicators, so a model enumerates $\sim$8--10 of them ---
making at least one refutable almost certain --- yet labels correctly, whereas
benign-looking phishing (wallet drainers, parked redirects) affords few claims and
is precisely where it mislabels. The contradicted-claim count therefore co-measures
how phish-rich a sample is, not how truthful the model is, and is unsafe as a
per-label quality gate.

\begin{table}[h]
\centering
\caption{Hallucination paradox: accuracy is higher on hallucinated samples than on
clean ones for all six providers. Llama is the only non-significant case
($p=0.10$).}
\label{tab:bench-halluc}
\small
\begin{tabular}{@{}lcccc@{}}
\toprule
Provider & acc$_{\text{hall}}$ & acc$_{\text{clean}}$ & $\Delta$ & OR \\
\midrule
GPT-5.1           & 0.898 & 0.650 & $+24.8$ & 4.72 \\
GPT-5.4           & 0.908 & 0.678 & $+23.0$ & 4.69 \\
DeepSeek-V4-Flash & 0.870 & 0.667 & $+20.4$ & 3.36 \\
Mistral-Large-3   & 0.931 & 0.742 & $+18.9$ & 4.69 \\
GPT-5.4-mini      & 0.900 & 0.725 & $+17.5$ & 3.41 \\
Llama-3.3-70B     & 0.899 & 0.834 & $+6.4$  & 1.76 \\
\bottomrule
\end{tabular}
\end{table}

\paragraph{The very-high-trust paradox.} Bucketing labels into trust quartiles,
accuracy climbs from Q1 to Q3 but then \emph{falls} in the top quartile Q4 --- for
every provider (Table~\ref{tab:bench-trust}; e.g.\ DeepSeek Q3 $0.95\!\to$ Q4
$0.69$). A near-saturated trust score requires many indicators, all verified or
unverifiable, which is easiest to achieve on samples with a clear \emph{surface}
match to the predicates --- a set that happens to include many benign pages with
login forms, hidden iframes, or redirects. The verifier audits claim-grounding,
not sample-level ambiguity, so a rule ``accept if $\tau\geq Q4$'' would
preferentially admit a worse-than-average slice. Both paradoxes point the same
way: trust belongs in the health monitor as a cross-round aggregate, never on an
individual label.

\begin{table}[h]
\centering
\caption{Trust-quartile accuracy. Q4 (highest trust) is below Q3 for every
provider.}
\label{tab:bench-trust}
\small
\begin{tabular}{@{}lcccc@{}}
\toprule
Provider & Q1 & Q2 & Q3 & Q4 \\
\midrule
DeepSeek-V4-Flash & 0.585 & 0.949 & 0.952 & 0.692 \\
GPT-5.4           & 0.554 & 0.942 & 0.975 & 0.785 \\
GPT-5.4-mini      & 0.587 & 0.950 & 0.959 & 0.798 \\
GPT-5.1           & 0.438 & 0.926 & 0.960 & 0.771 \\
Mistral-Large-3   & 0.766 & 0.940 & 0.968 & 0.765 \\
Llama-3.3-70B     & 0.654 & 0.976 & 0.916 & 0.895 \\
\bottomrule
\end{tabular}
\end{table}

\subsection{Reasoning styles across families}
\label{app:bench-style}

Providers differ in reasoning \emph{style} far more than in quality
(Table~\ref{tab:bench-style}), and the difference is legible in the verifier's
own resolution counts. Llama is \emph{concrete}: short, generic,
predicate-matchable indicators, which earns it the highest verify rate ($0.51$)
but the least forensic value. The GPT-5.4 family is \emph{abstract}: verbose
reasoning ($\sim$9.5 indicators per call) that cites HTML and OSINT fields by name
and is consequently classed \emph{unverifiable} most often ($0.58$). Two
consequences follow that bear directly on prompt design. First, only GPT-5.4 and
Llama actually exploit the prompt's OSINT dossier as named facts; DeepSeek
consumes it silently and Mistral cites ``domain age'' without ever quoting the
number. Second, GPT-5.4-mini pairs the highest self-confidence ($0.96$) with the
highest hallucination rate ($0.12$): its most decisive calls are also its most
fabricated --- something the verifier's trust captures (it lands mid-tier) but the
model's own confidence does not. The upshot is that a lower hallucination rate or
a higher verify rate is partly just a \emph{less specific} claim style, so neither
is a truthfulness ranking --- only a verbosity signature.

\begin{table}[h]
\centering
\caption{Reasoning-style profile. The verify/unverifiable split is a verbosity
signature, not a quality ranking; confidence and hallucination are not
self-correcting.}
\label{tab:bench-style}
\small
\begin{tabular}{@{}lcccc@{}}
\toprule
Provider & conf & verify & unver & contr \\
\midrule
Llama-3.3-70B     & 0.92 & \textbf{0.513} & 0.370 & 0.117 \\
GPT-5.4-mini      & 0.96 & 0.436 & 0.439 & \textbf{0.124} \\
GPT-5.1           & 0.91 & 0.395 & 0.514 & 0.091 \\
DeepSeek-V4-Flash & 0.96 & 0.357 & 0.545 & 0.098 \\
Mistral-Large-3   & 0.97 & 0.356 & 0.550 & 0.094 \\
GPT-5.4           & 0.93 & 0.339 & \textbf{0.582} & 0.079 \\
\bottomrule
\end{tabular}
\end{table}

\subsection{Chain-of-thought case studies}
\label{app:bench-cot}

Inspecting cached reasoning on shared sample IDs makes the styles --- and the
failure modes --- concrete.

\begin{itemize}[leftmargin=1.2em,itemsep=2pt,topsep=2pt]
\item \textbf{Textbook agreement} (sample 933,
\texttt{trezorrloggin.blogspot.com}). Ground truth phishing; all six providers
correct. GPT-5.4 cites the literal lookalike string and the iframe origin; Llama
simply lists ``typosquatting domain, hidden iframe'' --- the easy branch of the
AL pool, where any reasonable oracle suffices.
\item \textbf{Unanimous against ground truth} (sample 1430,
\texttt{imdb.com/list/\dots}). Labelled phishing in PP, yet all six call it benign
($0.93$--$0.99$ confidence), correctly reading the IMDb branding, the Amazon
script chain, and the 30-year domain age. This is a probable \emph{dataset}
mislabel --- the canonical \texttt{cross\_model\_split} pattern of
Appendix~\ref{app:bench-anom}, not a model error.
\item \textbf{A true cross-model blind spot} (sample 1679,
\texttt{gestaodecadastrosatualizar.blogspot.com}). Ground truth phishing; four of
five anchor on the Blogspot $404$ payload (``no login form'') and call benign,
missing that the URL reads ``update account records'' in Portuguese. Only Mistral
catches it, via its ``free Blogspot subdomain commonly abused'' prior --- a
reminder that the oracle helps only when at least one provider's prior is right.
\item \textbf{Lucky hallucination} (sample 869, a Roblox/Robux scam). Mistral
labels phishing (correct) but four of its supporting claims land as
\emph{contradicted}; trust drops, the label stays right. For retraining the label
is all that matters --- the fabricated support is irrelevant, which is precisely
why filtering on hallucination would have discarded a correct label.
\end{itemize}

\subsection{Anomaly catalogue and the dataset noise floor}
\label{app:bench-anom}

A retrospective extractor sorted $994$ audited verdicts into six operational
error classes (Table~\ref{tab:bench-anom}), and the two \emph{empty} classes are
as informative as the populated ones. \texttt{groundless\_phish}$=0$ across
$\sim$2{,}800 rows means every phishing verdict carries a named, checkable
indicator --- so the ``the LLM just post-hoc rationalises'' critique simply does
not apply to this oracle --- and \texttt{format\_failure}$=0$ confirms the JSON
contract holds for the healthy providers. The most consequential populated class
is \texttt{cross\_model\_split} ($243$): unanimous-LLM verdicts that contradict
the ground truth, on manual triage an estimated $70$--$80\%$ \emph{dataset}
mislabels --- legitimate IMDb lists, the YouTube homepage, German classifieds, all
labelled phishing in PP yet unanimously benign with strong OSINT legitimacy. This
is a genuine benchmark label-noise floor: it means part of the gap between the
LLM-AL result ($0.80$) and the ground-truth ceiling ($0.87$) is the dataset's own
noise, not the oracle's, and the true oracle ceiling is therefore higher than the
nominal numbers suggest.

\begin{table}[h]
\centering
\caption{Anomaly catalogue ($994$ verdicts). Two empty classes validate the
prompt; \texttt{cross\_model\_split} exposes a dataset label-noise floor.}
\label{tab:bench-anom}
\small
\begin{tabular}{@{}lrl@{}}
\toprule
Class & count & meaning \\
\midrule
inverse\_conservatism & 392 & benign call, evidence says phish \\
confident\_wrong      & 310 & high-confidence wrong verdict \\
cross\_model\_split   & 243 & unanimous LLMs vs.\ GT \\
lucky\_halluc         & 49  & correct label, refuted support \\
groundless\_phish     & \textbf{0} & phish with no named indicator \\
format\_failure       & \textbf{0} & unparseable output (healthy) \\
\bottomrule
\end{tabular}
\end{table}

\subsection{Health-monitor validation on real broken oracles}
\label{app:bench-llama}

The benchmark also stress-tests the health monitor of
Section~\ref{ssec:safety-layer} on the three oracles that failed for real, and
each isolates a different point. \textbf{Grok-4} returned \emph{zero} usable
labels in its first two rounds ($0/20$), stalled at the round-0 baseline, and is
excluded from the audit. \textbf{Kimi-K2.6} is the case that justifies checking
health at the \emph{parse} layer rather than the network layer: in a prior full
run all $1{,}232$ of its API calls returned HTTP~200 with zero transport errors,
yet $839$ ($68\%$) carried a body that did not parse to the schema, so only $33\%$
of queried labels were usable and the run recovered nothing ($+0.05$ points). A
network-level liveness check would have called the oracle perfectly healthy; the
monitor's parse-rate signal sees $0.33 \ll 0.80$ and halts within the first
rounds.

\textbf{Llama-3.3-70B} is the true-positive halt, and the cleanest demonstration
that the monitor is \emph{reactive}, not predictive (Figure~\ref{fig:llama-halt}).
Its parse rate is chronically marginal --- dipping repeatedly into the warn zone
and once critical ($0.75$ at round~5) --- and it finally strings together three
consecutive sub-$0.95$ rounds ($0.85$ and $0.80$ at rounds 16--17, and a third
below $0.95$ at round~18), tripping the sustained-parse rule and halting at
round~18. The monitor does not anticipate the collapse; it fires only \emph{after}
the malformed-JSON rate has stayed degraded for three rounds running. That the
downstream PP accuracy still looks healthy at the halt is both irrelevant to the
decision and invisible to it --- the monitor is label-free and watches the output
stream, never $F_{1,w}$. The halt nonetheless strictly improves the outcome
(Table~\ref{tab:llama}): running all $25$ rounds with the monitor disabled ends at
$\mathrm{PP_{test}}\,F_{1,w}=0.6115$, \emph{below} the round-0 baseline of
$0.6234$, whereas the monitored run stops at $0.6265$. Having flagged three broken
oracles while never tripping on the four healthy providers across $25$ rounds
each, the monitor is calibrated tightly enough to act and loosely enough to avoid
false alarms.

\begin{table}[h]
\centering
\caption{The monitored halt on Llama-3.3-70B is a true positive: disabling the
monitor and running to completion ends below the round-0 baseline ($0.6234$).}
\label{tab:llama}
\small
\begin{tabular}{@{}lcc@{}}
\toprule
Llama-3.3-70B run & rounds & $\mathrm{PP_{test}}\,F_{1,w}$ \\
\midrule
monitor on (halts) & 18 & \textbf{0.6265} \\
monitor off (full) & 25 & 0.6115 \\
\bottomrule
\end{tabular}
\end{table}

\subsection{Matched-budget provider cross-check}
\label{app:bench-smallbench}

Finally, to confirm that \emph{absolute} recovery --- not just per-label accuracy
--- is provider-insensitive, the loop is run end-to-end on a reduced PP subset
($|\mathrm{PP_{stream}}|{=}500$, $|\mathrm{PP_{test}}|{=}250$) at matched budget
across four providers and two selectors (Table~\ref{tab:smallbench}). The absolute
numbers are deliberately not comparable to the full-PP main results --- the
smaller stream and label budget lower the ceiling, and $\mathrm{PP_{test}}\,n{=}250$
carries $\mathrm{SE}\approx\pm2.5$~pp on $F_{1,w}$ --- so the table supports only a
\emph{relative} reading. On that reading, providers cluster within
$\approx 3$--$4$ points per selector, consistent with the benign-class-only
variance of Appendix~\ref{app:bench-acc}: under a fixed pipeline, the oracle is
interchangeable to within a few points, and the structured prompt is doing the
work.

\begin{table}[h]
\centering
\caption{Matched-budget cross-check ($\mathrm{PP_{stream}}{=}500$, seed $2025$).
Providers cluster within $\approx 4$ points per selector; absolute values are not
comparable to the full-PP main results.}
\label{tab:smallbench}
\small
\begin{tabular}{@{}llcc@{}}
\toprule
Selector & Provider & PP $F_{1,w}$ & PP $R_{\text{phish}}$ \\
\midrule
\multirow{4}{*}{drift-anchored}
 & DeepSeek-V4-Flash & 0.675 & 0.402 \\
 & Mistral-Large-3   & 0.671 & 0.418 \\
 & GPT-5.4           & 0.635 & 0.352 \\
 & GPT-5.4-mini      & 0.635 & 0.352 \\
\midrule
\multirow{4}{*}{pure-random}
 & GPT-5.4-mini      & 0.734 & 0.516 \\
 & Mistral-Large-3   & 0.717 & 0.500 \\
 & DeepSeek-V4-Flash & 0.701 & 0.467 \\
 & GPT-5.4           & 0.701 & 0.467 \\
\bottomrule
\end{tabular}
\end{table}

Taken together, the benchmark establishes the two facts the main paper needs from
it: the self-repair loop is robust to the \emph{choice} of oracle --- accuracy
spans a modest $8.9$ points across six LLMs, all of the variance on the benign
class, led by open-weights and European models --- and it is safe under oracle
\emph{failure}, with the verifier exposing two paradoxes that confirm its trust is
observability rather than a gate, and the health monitor validated on three real
broken oracles including a true-positive halt that strictly improves the outcome.
          % Appendix

% --- Bibliography ----------------------------------------------------
\bibliographystyle{ACM-Reference-Format}
\bibliography{bibliography}

\end{document}
